\NativeSlot{61}{156}{%
\ExamGroup{四、操作。{\normalfont\kaishu\fontsize{10.8bp}{14bp}\selectfont （共18分）}}

% p039_s01

\ExamQuestion{1.}{5}%
按要求，画一画。（每两个格点之间的距离为1厘米）\par\noindent （1）以图中线段为三角形的底边，画一个等腰三角形，并画出底边上的高。\par
\BookFigure{figures/p039_s01-01.tex}

% p039_s02

\ExamQuestion{（2）}{5}%
画一个梯形，并把梯形分成一个平行四边形和一个三角形。\par
\BookFigure{figures/p039_s02-01.tex}
}
\NativeSlot{336}{81}{%
% p039_s03

\ExamQuestion{2.}{5}%
画一画，写一写。\par\noindent （1）利用直尺和圆规，画出以线段\(a\)为底边，线段\(b\)和线段\(c\)为边的三角形。（保留作图痕迹）\par
\BookFigure{figures/p039_s03-01.tex}
}
\NativeSlot{490}{36}{%
% p039_s04

\ExamQuestion{（2）}{7}%
根据你画的三角形，请提出一个值得思考的数学问题。（不用解答）\par
}
\NativeSlot{565}{130}{%
\ExamGroup{五、解决问题。{\normalfont\kaishu\fontsize{10.8bp}{14bp}\selectfont （共22分）}}

% p039_s05

\ExamQuestion{1.}{7}%
我国古代很早就用方正图形做器物，讲究几何之美。手工课上，李明用一根铁丝先围成了边长9厘米的正方形，如果要用这块铁丝围成一个等边三角形，它的边长是多少厘米？（7分）\par
}
